\section{The Channels Together}
\label{sec:together}

Figure~\ref{fig:alone} shows four channels, each on its own, with the same outcomes and
the same scale. In the Philippine example, electricity costs have by far the largest effect
on the economy. Public investment and the household carbon tax move the economy little.
The change in generation technology reshapes the electricity sector, as described above,
but has little effect on GDP.

\begin{figure}[H]
\centering
\includegraphics[width=\linewidth]{channels-alone.pdf}
\caption{Each channel on its own: GDP, consumption, capital and labour, percent change from
the control, Philippine example. All panels share the same scale.}
\label{fig:alone}
\end{figure}

The Philippine scenario combines electricity costs, public infrastructure and health. The
carbon price is passed to CLEWS for its next solution rather than applied in the economy,
and the scenario does not include a change in generation technology or investment
incentives. Without health, the combined result is essentially the effect of electricity
costs: GDP is about 0.19 percent lower over the first decade and 0.14 percent lower in 2075,
and consumption is about 0.07 and 0.11 percent lower, respectively. Public investment adds
almost nothing in this scenario.

\pending{the combined result with health, over time, from the corrected health run.}

The distribution of these effects is where the overlapping-generations structure matters
most. OG-Core follows how the change in prices, wages, returns and taxes affects households
of each age and lifetime income, and how the same scenario affects those who are working,
those who are retired, and those not yet born.

\pending{who bears the cost: change in lifetime welfare by age and lifetime income, from the
corrected combined run.}
